% =====================================================================
% Section 2 --- From transformation research to input-output models
% ---------------------------------------------------------------------
% Job: why *plural* IO models are the bridge; the construction and price
%      basis of the empirically informed green-investment vector; light
%      mathematics, the beginner on-ramp.
% Mindmap: N08, N09, N10.
% Status: STUB --- narrative outline v3 (2026-09-20).
% =====================================================================
% NOTE (numbers belong here, not in the abstract) -- per ADR-0024:
%   G0 = 40.3 bn in 2019 prices = the horizon mean (2024-2050) = 1.331 % of GDP
%        (40.3 / 3027.8; equals the booked B_gov = 0.01330991);
%   58.9 bn/yr in 2023 prices = the raw impulse mean (the CSV's own basis);
%   2024 first-year peak 81 bn (2023 prices) ~ 1.9 % of GDP;
%   deflator 1.46 (construction-cost indices); 58.89 / 1.46 = 40.338 = 40.337 bn.
% Cite BOTH price bases; never mix them in one sentence (outline external
% review, "The GDP share of the programme").
\section{From transformation research to input--output models}\label{sec:io-bridge}

% TODO (drafting notes):
% - Open concrete: the npj Climate Action impulse (Hornykewycz et al. 2025),
%   ~58 bn/yr and 3.1 tn until 2050 in 2023 prices; the 2024 first-year peak
%   ~81 bn ~ 1.9 % of GDP; G0 = 40.3 bn in 2019 prices = the horizon mean,
%   1.331 % of GDP (ADR-0024: both price bases cited, never mixed).
% - Why sectoral IO models are the applied lingua franca of transformation
%   research (Creutzig et al. 2018; Miller & Blair; Leontief 1936, 1986).
% - The vector: `impulses.csv`, incidence psi, the deflator 1.46 and the
%   ADR-0024 price-basis verification (defer the full provenance to App. D).
% - The beginner on-ramp: keep the algebra light here; the formal kernel is
%   Section 4 and Appendix B.

% ===== CARRIED OVER FROM AN OLD DRAFT =============================
% [devised for an old draft; not in the validated submission]
% source: revised_manuscript/chapters/ch04_data.tex (commit 41144dd)
% =================================================================
% To provide a concrete and reliable empirical example for such investment requirements, we take the German residential housing sector as our main case of analysis. When doing so we base our analysis and discussion on a specific estimate of the costs associated with a transformation of the German housing sector towards carbon neutrality as supplied by \citet{Hornykewycz.2025}. This study estimates expected costs by proportionally scaling past renovations costs under the assumption that an increase in the renovation intensity takes place, that is sufficient to meet Germany's climate goals. Specifically, the study at hand finds that renovation efforts will have to be more than doubled, if Germany is set to comply with its overall goal to achieve carbon neutrality till 2050. Such a shift would amount to yearly costs of about 40.3 bn. € in 2019 prices\footnote{We assumed a deflator of 1.46 based on price indices for construction costs as specified in \citet{Hornykewycz.2025}. In 2023 prices the shock amounts to about 58 bn. €.}, which would represent a substantial shift in final demand. As the study also details how these costs should be mapped on different sectors, it provides sufficient information to construct a concrete and directly applicable specification for the sector-specific annual changes in final demand imposed by such a reorientation. This specification shocks final demand in seven sectors, where sectors are sorted by the relative increase associated with these shocks (as given in brackets): 
% Glass and glassware (+50.3\%), ceramic products and processed stones (+37,3 \%), specialized construction work (+27,8\%), rubber and plastic products (+13.4\%), chemicals and chemical products (+3.7\%), machinery (+1.2\%) and electrical equipment (+1.0\%).
% In terms of the data used, we apply the so estimated shock vector as proposed by \citet{Hornykewycz.2025} with official German input-output data for 2019, the most recent year available \citep[see][]{destatis_strukturdaten}.

